\section{代数方程的根式可解性}\label{sec:galois-2-8}
正如本章引言中已经提到的，Galois 理论能够解决代数方程的根式可解性问题. 下面我们说明这一点.
\subsection{多项式的 Galois 群}\label{sec:galois-2-8-1}
设 $f\in k[X]$ 是可分多项式. 称 $f$ 在 $k$ 上的一个分裂域 $K_f$ 的 Galois 群 $G_f$ 为 $f$ 的 Galois 群. 这个群置换 $f$ 的根，其在 $K_f$ 上的作用由在这些根上的作用决定，因为这些根生成 $K_f$. 因此，$G_f$ 可视为 $f$ 的根的置换群的一个子群. 若将 $f$ 在 $K_f$ 中的根编号为 $\alpha_1,\ldots,\alpha_n$，则 $G_f$ 可视为 $\mathfrak S_n$ 的一个子群.
\begin{lemma*}
$f$ 在 $k[X]$ 中不可约，当且仅当 $G_f$ 传递地置换 $f$ 在 $K_f$ 中的根.
\end{lemma*}
\begin{proof}
我们已经多次见过这一点. 再重复一下论证. 若 $f$ 不可约，且 $\alpha,\beta$ 是它的两个根，则存在唯一一个把 $\alpha$ 映为 $\beta$ 的 $k$-同构 $k[\alpha]\xrightarrow\sim k[\beta]$；由于 $K_f$ 正规，此同构延拓为自同构 $K_f\xrightarrow\sim K_f$. 反之，Galois 扩张的性质 ii) 告诉我们，若 $G_f$ 在 $f$ 的根上传递地作用，则 $f$ 是其每个根的最小多项式，特别地不可约.
\end{proof}
若反而有分解 $f=f_1f_2\in k[X]$，令 $K_{f_i}$ 为 $K_f$ 中由 $f_i$ 的根生成的子域. 由于 $K_{f_i}$ 在 $k$ 上是 Galois 扩张，存在满态射 $G_f\twoheadrightarrow G_{f_i}$，其核为 $\operatorname{Gal}(K_f/K_{f_i})$. 乘积态射 $G_f\longrightarrow G_{f_1}\times G_{f_2}$ 则是单射，因为 $K_f$ 由 $K_{f_1}.K_{f_2}$ 生成.
\subsection{根式可解性与可解群}\label{sec:galois-2-8-2}
回顾，多项式 $f\in\mathbb Q[X]$ 称为``根式可解的''，若它的根能从有理数出发，依次施行 $+,-,\cdot,\div$ 和 $\sqrt[n]{x}$ 中的运算表示出来.
\begin{definition*}
有限群 $G$ 称为\newterm{galois-solvable-group}{可解群}{solvable group}，若它有递减的正规子群列 $G=G_0\supset G_1\supset\cdots\supset G_r=\{1\}$，使 $G_i/G_{i+1}$ 是 Abel 群.
\end{definition*}
\begin{exercise*}
设 $H$ 是 $G$ 的子群. 证明：
\begin{itemize}
\item $G$ 可解 $\Rightarrow H$ 可解.
\item 若 $H$ 正规，则 $G$ 可解 $\Leftrightarrow H$ 与 $G/H$ 均可解.
\end{itemize}
\end{exercise*}
Galois 定理可以表述如下：
\begin{theorem*}
多项式 $f$ 根式可解，当且仅当它的 Galois 群可解.
\end{theorem*}
\begin{proof}
先设 $f$ 根式可解. 这等价于 $K_f$ 包含在扩张塔 $\mathbb Q=K_0\subset K_1\subset\cdots\subset K_r$ 的``最高一层'' $K_r$ 中，其中对每个 $i=1,\ldots,r$，都有 $K_i=K_{i-1}(\alpha_i)$，且存在 $n_i\in\mathbb N$ 使 $\alpha_i^{n_i}\in K_{i-1}$. 可以选取此塔，使 $K_1$ 是 $n$ 次分圆扩张，其中 $n$ 为各 $n_i$ 的最小公倍数. 根据\ref{sec:galois-2-7-5} 和\ref{sec:galois-2-7-4}，每个 $K_i\supset K_{i-1}$ 都是 Galois 群为 Abel 群的 Galois 扩张. 遗憾的是，$i>2$ 时 $K_i$ 未必在 $\mathbb Q$ 上是 Galois 扩张. 因此，在 $\overline{\mathbb Q}$ 中用其 Galois 闭包 $K_i'$ 代替 $K_i$. 于是得到 Galois 扩张塔 $\mathbb Q=K_0\subset K_1=K_1'\subset K_2'\subset\cdots\subset K_r'$，其中
\begin{align*}
K_i'=K_{i-1}'(\alpha_i^{(j)},\ j=1,\ldots,r_i),
\end{align*}
这里 $\alpha_i^{(j)}$ 表示 $\alpha_i$ 在 $\overline{\mathbb Q}$ 中的各个共轭元. 此时 $(\alpha_i^{(j)})^{n_i}$ 是 $\alpha_i^{n_i}$ 的一个共轭元，因而属于 $K_{i-1}'$，由\ref{sec:galois-2-7-5} 的推论，$\operatorname{Gal}(K_i'/K_{i-1}')$ 为 Abel 群.

用 Galois 对应将上述事实翻译过来. 记 $G_i':=\operatorname{Gal}(K_r'/K_i')$，它是 $G_0'=\operatorname{Gal}(K_r'/\mathbb Q)$ 的子群. 各 $G_i'$ 都在 $G_0'$ 中正规，且相继商群 $G_i'/G_{i+1}'$ 都是 Abel 群. 因此 $G_0'$ 可解. 由于 $K_f\subset K_r'$，群 $G_f:=\operatorname{Gal}(K_f/\mathbb Q)$ 是 $G_0$ 的商群，故也可解. 事实上，若 $G_{f,i}$ 是 $G_i'$ 在 $G_f$ 中的像，则每个 $G_{f,i}$ 都在 $G_f$ 中正规，且相继商群 $G_{f,i}/G_{f,i+1}$ 是 Abel 群，因为它们是 $G_i/G_{i+1}$ 的商群.\footnote{译注：本段原文两处省略撇号，分别写作 $G_0$ 和 $G_i/G_{i+1}$；按上下文应为 $G_0'$ 与 $G_i'/G_{i+1}'$.}

反过来，现在设 $G_f=\operatorname{Gal}(K_f/\mathbb Q)$ 可解. 记 $K_f'$ 为由 $K_f$ 和 $n$ 次单位根生成的域，其中 $n=[K_f:\mathbb Q]$. 它也是 $\mathbb Q$ 的 Galois 扩张，其 Galois 群 $G_f'$ 满射到 $G_f$，核 $\operatorname{Gal}(K_f'/K_f)$ 为 Abel 群（其阶整除 $\varphi(n)$）. 因此 $G_f'$ 也可解. 记 $G_{f,1}':=\operatorname{Gal}(K_f'/\mathbb Q(e^{2i\pi/n}))$；它是 $G_f'$ 的正规子群，商群 $G_f'/G_{f,1}'$ 为 Abel 群（同构于 $(\mathbb Z/n\mathbb Z)^\times$）. 由于 $G_f'$ 可解，存在正规子群 $G_{f,1}'\supset G_{f,2}'\supset\cdots\supset G_{f,r}'=\{1\}$，使 $G_{f,i}'/G_{f,i+1}'$ 为 Abel 群. 事实上，每个有限 Abel 群均为循环群的积，故甚至可以设 $G_{f,i}'/G_{f,i+1}'$ 为循环群. 注意，对 $i\geq1$，$G_{f,i}'/G_{f,i+1}'$ 的阶整除 $G_f'/G_{f,1}'$ 的阶，后者等于次数 $[K_f':\mathbb Q(e^{2i\pi/n})]$，而这个次数整除 $[K_f:\mathbb Q]=n$.\footnote{译注：原文此处把 $G_f'/G_{f,1}'$ 的阶与固定域之上的次数相等同；该次数实际是 $|G_{f,1}'|$. 此外，循环商群的细分仅需逐级正规，不能一概要求均在全群中正规. 原文表述保留.}
令 $K_{f,i}':=(K_f')^{G_i}$.\footnote{译注：原文此处写 $G_i$，上下文对应的群为 $G_{f,i}'$.} Galois 对应告诉我们，塔
\begin{align*}
\mathbb Q\subset\mathbb Q(e^{2i\pi/n})=K_{f,1}'\subset K_{f,2}'\subset\cdots\subset K_{f,r}'=K_f'
\end{align*}
由 Galois 扩张组成，其中 $K_{f,i}'/K_{f,i+1}'$ 的 Galois 群为阶 $n_i$ 整除 $n$ 的循环群. 根据\ref{sec:galois-2-7-5} 的定理，这样的扩张形如 $K_{f,i}'=K_{f,i+1}'(\sqrt[n_i]{a_i})$.\footnote{译注：原文这两处扩张方向与所列递增塔相反，应涉及 $K_{f,i+1}'/K_{f,i}'$ 和 $K_{f,i+1}'=K_{f,i}'(\sqrt[n_i]{a_i})$. 此处保留原式并指出其方向问题.} 因而 $f$ 根式可解.
\end{proof}
\subsection{次数不超过四的方程的可解性}\label{sec:galois-2-8-3}
在 Galois 的时代，人们早已知道次数不超过四的多项式根式可解. 下面给出一个概念性的解释. 令 $n:=\deg(f)$. 我们已经看到，$G_f:=\operatorname{Gal}(K_f/K)$ 在 $f$ 的根集上的置换作用给出嵌入 $G_f\hookrightarrow\mathfrak S_n$（将根编号之后）.\footnote{译注：原文此处基域记作 $K$，而前后讨论的基域为 $\mathbb Q$.} 当 $n\leq4$ 时，群 $\mathfrak S_n$ 可解. 事实上，$\mathfrak S_2=\mathbb Z/2\mathbb Z$ 为 Abel 群；$\mathfrak S_3$ 通过符号满射到 $\mathbb Z/2\mathbb Z$，核为 $\mathfrak A_3=\mathbb Z/3\mathbb Z$；$\mathfrak S_4$ 满射到 $\mathbb Z/2\mathbb Z$，核为 $\mathfrak A_4$，后者的子群 $\{\mathrm{id},(1,2)(3,4),(1,3)(2,4),(1,4)(2,3)\}$ 同构于 $(\mathbb Z/2\mathbb Z)^2$，是正规子群，且商群同构于 $\mathbb Z/3\mathbb Z$. 由于可解群的每个子群均可解（练习），只要 $\deg(f)\leq4$，$G_f$ 便可解.
\subsection{一个五次方程的不可解性}\label{sec:galois-2-8-4}
第一个根式不可解的代数方程例子归功于 Abel. 不过，Galois 理论为 Abel 的例子提供了更具概念性的解释.
\begin{lemma*}
$n\geq5$ 时，群 $\mathfrak S_n$ 不可解.
\end{lemma*}
\begin{proof}
只需证明 $\mathfrak A_n$ 不可解. 为此，只需证明 $\mathfrak A_n$ 没有 Abel 商群. 这等价于证明，由 $\mathfrak A_n$ 中元素的换位子 $xyx^{-1}y^{-1}$ 生成的子群 $[\mathfrak A_n,\mathfrak A_n]$ 等于 $\mathfrak A_n$. 事实上，每个从 $\mathfrak A_n$ 到 Abel 群 $G$ 的态射都在 $[\mathfrak A_n,\mathfrak A_n]$ 上平凡.\footnote{译注：``没有 Abel 商群''应理解为``没有非平凡 Abel 商群''.}

回顾，$\mathfrak A_n$ 由三循环生成. 事实上，只需证明两个换位的积 $\tau=(i,j)(k,l)$ 是三循环的积. 若 $\{i,j\}=\{k,l\}$，则 $\tau=\mathrm{id}$；若 $|\{i,j\}\cap\{k,l\}|=1$，例如设 $j=k$，则 $\tau=(i,j,l)$；若 $\{i,j\}\cap\{k,l\}=\varnothing$，则 $\tau=(i,j)(j,k)(j,k)(k,l)=(i,j,k)(j,k,l)$.

所以只需证明每个三循环都是 $\mathfrak A_n$ 中的换位子. 公式
\begin{align*}(i,j,k)=(i,j)(j,k)=(i,j)(i,k)(i,j)^{-1}(i,k)^{-1}\end{align*}
表明 $(i,j,k)$ 是 $\mathfrak S_n$ 中的换位子. 为把它变成 $\mathfrak A_n$ 中的换位子，选取不同于 $i,j,k$ 的 $l\ne m$；由于 $n\geq5$，这是可能的. 此时 $(l,m)$ 与 $(i,j)$ 和 $(i,k)$ 都交换. 因而令 $\tau=(i,j)(l,m)$ 和 $\sigma=(i,k)(l,m)$，便有 $\tau\sigma\tau^{-1}\sigma^{-1}=(i,j,k)$，且 $\tau,\sigma\in\mathfrak A_n$.
\end{proof}
\begin{remark*}
事实上，更强的结论成立：当 $n\geq5$ 时，群 $\mathfrak A_n$ 是单群，即没有真且非平凡的正规子群.
\end{remark*}
现在我们要构造一个 Galois 群为 $\mathfrak S_5$ 的五次多项式. 为此，下述引理会有帮助.
\begin{lemma*}
群 $\mathfrak S_n$ 由任意一对由一个 $n$ 循环和一个换位组成的元素 $(\sigma,\tau)$ 生成.
\end{lemma*}
\begin{proof}
设 $\tau=(i,j)$. 必要时以 $\sigma$ 的某个幂代替 $\sigma$，可以假设 $j=\sigma(i)$. 于是，对每个 $s=0,\ldots,n-1$，有 $\sigma^s\tau\sigma^{-s}=(\sigma^s(i),\sigma^{s+1}(i))$. 设 $r>s$. 则
\begin{align*}
&(\sigma^{r-1}(i),\sigma^r(i))\cdots(\sigma^{s+1}(i),\sigma^{s+2}(i))(\sigma^s(i),\sigma^{s+1}(i))\\
&\quad\cdot(\sigma^{s+1}(i),\sigma^{s+2}(i))\cdots(\sigma^{r-1}(i),\sigma^r(i))
=(\sigma^s(i),\sigma^r(i)).
\end{align*}
这表明 $\sigma,\tau$ 生成的子群包含所有换位，因此等于 $\mathfrak S_n$.
\end{proof}
\noindent\textbf{译注.} 原引理对一般 $n$ 不成立，例如 $(1,2,3,4)$ 与 $(1,3)$ 不能生成 $\mathfrak S_4$. 证明中以幂代替循环的步骤需要相应指数与 $n$ 互素. 对下文所用的素数 $n=5$，这一条件成立.

因此，我们希望找到一个五次多项式，使其 Galois 群包含一个五循环和一个换位. 注意到：
\begin{lemma*}
设 $f\in k[X]$ 可分. 若 $f$ 不可约且次数为素数 $p$，则 $G_f$ 在 $f$ 的根的置换群中包含一个 $p$ 循环.
\end{lemma*}
\begin{proof}
$f$ 的每个根域的次数均为 $p$（由于 $f$ 不可约，它同构于 $k[X]/(f)$），所以 $p\mid[K_f:k]$，因而 $G_f$ 包含一个 $p$ 阶元素. 但 $\mathfrak S_p$ 中的 $p$ 阶元素只有 $p$ 循环.
\end{proof}
寻找不可约多项式时，下述判别法非常有用.
\begin{proposition*}[{Eisenstein 判别法}]
设 $f=X^n+a_{n-1}X^{n-1}+\cdots+a_0\in\mathbb Z[X]$. 假设存在素数 $p$，对每个 $i$ 都有 $p\mid a_i$，但 $p^2\nmid a_0$. 则 $f$ 在 $\mathbb Q[X]$ 中不可约.
\end{proposition*}
\begin{proof}
设 $f=gh$ 是 $\mathbb Q[X]$ 中的分解，$g,h$ 首一. 我们已经说明此时 $g,h\in\mathbb Z[X]$. 写成 $g=X^m+b_{m-1}X^{m-1}+\cdots+b_0$ 和 $h=X^r+c_{r-1}X^{r-1}+\cdots+c_0$. 由 $b_0c_0=a_0$，$p$ 不整除 $b_0$ 或不整除 $c_0$. 设 $p$ 不整除 $c_0$，令 $k$ 为使 $p$ 不整除 $b_k$ 的最小整数（由于 $b_m=1$，它确实存在）. 等式 $a_k=\sum_{i+j=k}b_ic_j$ 表明 $p$ 不整除 $a_k$，故 $k=n$，$h(X)=1$.
\end{proof}
现在希望找到一种方法，在 $\mathbb Q[X]$ 中不可约多项式的 Galois 群中构造换位. 诀窍是利用复共轭，它至少给出一个二阶自同构. 事实上，设 $f$ 恰有两个根位于 $\mathbb C\setminus\mathbb R$ 中，则复共轭交换这两个根，固定其余所有根，从而给出 $G_f$ 中的一个换位.
\begin{corollary*}
多项式 $f=X^5-10X+5\in\mathbb Q[X]$ 的 Galois 群是 $\mathfrak S_5$. 特别地，$f$ 根式不可解.
\end{corollary*}
\begin{proof}
取 $p=5$，Eisenstein 判别法表明 $f$ 在 $\mathbb Q[X]$ 中不可约，故由最后一个引理，$G_f$ 包含五循环. 另一方面，函数 $\mathbb R\longrightarrow\mathbb R,\ x\longmapsto f(x)$ 的单调性表表明，$f$ 有三个实根，所以 $G_f$ 包含一个换位. 因而 $G_f\simeq\mathfrak S_5$.
\end{proof}

